\begin{proofbox}
    \textbf{\textcolor{CProof}{证明目标与前提}}
    设 $n\ge1$、$a_n\ne0$、$a_0\ne0$，且 $\dfrac pq$ 为 $f(x)=0$ 的有理根，其中 $p\ne0$、$q>0$、$p,q$ 互质。证明
    \[
        p\mid a_0,\qquad q\mid a_n.
    \]
    前三法着重证明有理根定理；第四法进一步严格证明试根后可以提出整数系数的一次因式。

    \medskip
    \textbf{\textcolor{CProof}{方法一（门槛最低：清分母，逐项观察）}}
\begin{center}
\begin{tikzpicture}[x=1cm,y=1cm,font=\small]
\node[font=\bfseries] at (5.5,3.6) {清分母后：每一列是一项（以三次为例）};
\foreach \x/\t in {2.4/{$a_3p^3$},4.8/{$a_2p^2q$},7.2/{$a_1pq^2$},9.6/{$a_0q^3$}} {
\node at (\x,2.65) {\t};
}
\node at (0.55,1.65) {含 $p$};
\node at (0.55,0.65) {含 $q$};
\foreach \x in {2.4,4.8,7.2,9.6} {
\draw[gray!50] (\x-0.8,0.2) rectangle (\x+0.8,2.1);
}
\foreach \x in {2.4,4.8,7.2} {\fill[CStatement] (\x,1.65) circle (0.13);}
\foreach \x in {4.8,7.2,9.6} {\fill[CExample] (\x,0.65) circle (0.13);}
\node at (9.6,1.65) {未必};\node at (2.4,0.65) {未必};
\draw[CStatement,thick] (1.6,2.2)--(8,2.2);
\draw[CExample,thick] (4,-0.05)--(10.4,-0.05);
\node at (5.5,-0.65) {总和为零：另三项有公因子，剩下一项也必须有。};
\end{tikzpicture}
\end{center}


    将 $f(p/q)=0$ 两边同乘 $q^n$，得整数等式
    \[
      a_np^n+a_{n-1}p^{n-1}q+\cdots+a_1pq^{n-1}+a_0q^n=0.\tag{1}
    \]
    图中把三次情形各项包含的公因子逐列标出；一般 $n$ 次时规律相同。
    \begin{enumerate}[leftmargin=*,itemsep=0.6em]
    \item \textbf{看分子：} 除最后一项外均含 $p$。移项后得到 $p\mid a_0q^n$。因 $p,q$ 互质，$|p|$ 的全部质因子及相应次数只能由 $a_0$ 提供，故 $p\mid a_0$。若 $|p|=1$，该结论直接成立。
    \item \textbf{看分母：} 除第一项外均含 $q$。同理 $q\mid a_np^n$；由 $q,p^n$ 互质，得 $q\mid a_n$。若 $q=1$，该结论直接成立。
    \end{enumerate}
    因此 $\boxed{p\mid a_0,\ q\mid a_n}$。负分子的整除关系按绝对值判断；\textbf{互质}是去掉 $q^n$ 或 $p^n$ 的关键。
\end{proofbox}
\clearpage
\begin{proofbox}
    \textbf{\textcolor{CProof}{方法二（速度最快：两端移项法）}}
\begin{center}
\begin{tikzpicture}[x=0.75cm,y=0.75cm,font=\small]
\node[font=\bfseries] at (7,5.5) {$3x^2-x-2=0,\quad p=-2,\ q=3$：清分母得 $12+6-18=0$};
\node[anchor=west] at (0,4.6) {按 $q=3$ 分组：$12=18-6$};
\foreach \x in {0,...,5} {\foreach \y in {0,1,2} {
\filldraw[fill=CExample!15,draw=white] (\x,\y+1) rectangle +(1,1);
}}
\foreach \x in {0,...,3} {\foreach \y in {0,1,2} {
\filldraw[fill=CStatement!40,draw=white] (\x,\y+1) rectangle +(1,1);
}}
\foreach \y in {1,2,3} {\draw[gray] (0,\y) rectangle (6,\y+1);}
\node at (2,0.35) {$12=3\times4$};\node at (5,0.35) {$6=3\times2$};
\node[anchor=west] at (8,4.6) {按 $|p|=2$ 分组：};
\foreach \x in {8,...,13} {\foreach \y in {0,1,2} {
\filldraw[fill=CStatement!25,draw=white] (\x,\y+1) rectangle +(1,1);
}}
\foreach \x in {8,10,12} {\foreach \y in {1,2,3} {\draw[gray] (\x,\y) rectangle +(2,1);}}
\node at (11,0.35) {$18=12+6=2\times9$};
\node at (7,-0.65) {每小格表示 $1$；分组没有余格，表示整除。};
\end{tikzpicture}
\end{center}

    \begin{description}[style=nextline, leftmargin=2.8em, labelsep=0.5em, itemsep=0.48em]
        \item[\textbf{步骤一（抓住整数等式）：}]
            由 $f(p/q)=0$ 清分母，立即写出
            \[
               a_np^n+qU=0,\qquad pV+a_0q^n=0,
            \]
            其中 $U,V$ 均为整数。\textbf{理由：} 第一式把除首项外的项合并；第二式把除末项外的项合并。
        \item[\textbf{步骤二（一步判定）：}]
            由 $a_np^n=-qU$ 得 $q\mid a_np^n$，因 $p,q$ 互质，得 $q\mid a_n$；由 $a_0q^n=-pV$ 得 $p\mid a_0q^n$，同理得 $p\mid a_0$。
            \[
               \boxed{q\mid a_n,\quad p\mid a_0.}
            \]
            \textbf{理由：} 省去逐项书写，直接利用两端的整除关系。
    \end{description}

    \medskip
\end{proofbox}
\clearpage
\begin{proofbox}
    \textbf{\textcolor{CProof}{方法三（步骤最少：已知整数系数一次因式时）}}
\begin{center}
\begin{tikzpicture}[x=1cm,y=1cm,font=\small]
\node[font=\bfseries] at (5.6,3.8) {逐格相乘：只有左上格、右下格贡献首尾项};
\foreach \x/\v in {3/{$b_2x^2$},6/{$b_1x$},9/{$b_0$}} {\node at (\x,2.85) {\v};}
\node at (0.6,1.85) {$qx$};\node at (0.6,0.65) {$-p$};
\fill[CStatement!15] (1.5,1.25) rectangle (4.5,2.45);
\fill[CExample!15] (7.5,0.05) rectangle (10.5,1.25);
\draw[gray] (1.5,0.05) rectangle (10.5,2.45);
\draw[gray] (1.5,1.25)--(10.5,1.25);
\foreach \x in {4.5,7.5} {\draw[gray] (\x,0.05)--(\x,2.45);}
\node[text=CStatement] at (3,1.85) {$qb_2x^3$};
\node at (6,1.85) {$qb_1x^2$};\node at (9,1.85) {$qb_0x$};
\node at (3,0.65) {$-pb_2x^2$};\node at (6,0.65) {$-pb_1x$};
\node[text=CExample] at (9,0.65) {$-pb_0$};
\node at (5.6,-0.65) {三次示例：$a_3=qb_2$，$a_0=-pb_0$；一般次数同理。};
\end{tikzpicture}
\end{center}

    \begin{description}[style=nextline, leftmargin=2.8em, labelsep=0.5em, itemsep=0.48em]
        \item[\textbf{步骤一（利用整数系数分解）：}]
            利用\textbf{方法四将证明}的事实：若 $p/q$ 是整数系数多项式的最简有理根，则存在整数 $b_0,\ldots,b_{n-1}$ 使
            \[
             f(x)=(qx-p)(b_{n-1}x^{n-1}+\cdots+b_0).
            \]
            \textbf{理由：} 此处把“商式是整数系数”当作已经证明的结论，不能仅由因式定理擅自假定。
        \item[\textbf{步骤二（只比首尾两项）：}]
            乘开后比较最高次项与常数项，得
            \[
               a_n=qb_{n-1},\qquad a_0=-pb_0.
            \]
            由于 $b_{n-1},b_0$ 是整数，立即得到
            \[
               \boxed{q\mid a_n,\qquad p\mid a_0.}
            \]
            \textbf{理由：} 两项系数直接呈现“整数倍”关系，证明最短。
    \end{description}

    \medskip
\end{proofbox}
\clearpage
\begin{proofbox}
    \textbf{\textcolor{CProof}{方法四（常规通用：证明整数系数因式确实存在）}}
\begin{center}
\begin{tikzpicture}[x=1cm,y=0.85cm,font=\small]
\node[font=\bfseries] at (5.6,5.8) {除以 $2x-1$：逐列消去最高次项};
\foreach \x/\v in {4/{$x^3$},6/{$x^2$},8/{$x$},10/{$1$}} {\node at (\x,4.9) {\v};}
\foreach \y/\lab in {4/{原式},3/{减去 $(2x-1)3x^2$},2/{减去 $(2x-1)5x$},1/{减去 $(2x-1)2$}} {\node[anchor=east] at (2.8,\y) {\lab};}
\foreach \x/\y/\v in {4/4/6,6/4/7,8/4/-1,10/4/-2,4/3/0,6/3/10,8/3/-1,10/3/-2,4/2/0,6/2/0,8/2/4,10/2/-2,4/1/0,6/1/0,8/1/0,10/1/0} {
\draw[gray!40] (\x-0.85,\y-0.42) rectangle (\x+0.85,\y+0.42);
\node at (\x,\y) {$\v$};
}
\foreach \x/\y in {4/4,6/3,8/2} {\draw[CStatement,very thick] (\x-0.65,\y-0.34)--(\x+0.65,\y-0.34);}
\node at (5.6,-0.1) {每行记录相减后的系数；商为 $3x^2+5x+2$，余数为 $0$。};
\end{tikzpicture}
\end{center}


    \textbf{图中每次消去一项，所有系数仍为整数。一般证明如下。}
    \begin{enumerate}[leftmargin=*,itemsep=0.65em]
    \item 对任意整数系数多项式 $P(x)=c_mx^m+\cdots+c_0$（$m\ge1$），若 $P(p/q)=0$，清分母得 $c_mp^m+qW=0$，其中 $W\in\mathbb Z$。由互质性得 $q\mid c_m$。
    \item 取整数 $b=c_m/q$，作差
    \[
      R(x)=P(x)-(qx-p)bx^{m-1}.
    \]
    因 $qb=c_m$，最高次项恰好消去。$R$ 仍为整数系数多项式，且 $R(p/q)=P(p/q)-0=0$。
    \item 从 $P=f$ 开始，若余下多项式次数仍至少为 $1$，重复前两步；每次次数严格下降，故有限次后只剩常数 $C$（也可能提前得到零多项式）。于是
    \[
      f(x)=(qx-p)(b_{n-1}x^{n-1}+\cdots+b_0)+C,
    \]
    所有 $b_j,C$ 都是整数；跳过的次数对应 $b_j=0$。
    \item 代入 $x=p/q$，得 $0=C$，所以确实存在整数系数商式。再比较首尾项：
    \[
      a_n=qb_{n-1},\qquad a_0=-pb_0
      \quad\Longrightarrow\quad q\mid a_n,\ p\mid a_0.
    \]
    \end{enumerate}
    本法独立使用清分母与互质性，不依赖方法三；因此也补足了方法三所需的整数系数前提。\textbf{只有验出根后才能提出 $qx-p$}，不能把候选值直接当作根。
\end{proofbox}
\clearpage
